\documentclass[french]{article}
\usepackage{graphicx} % Required for inserting images

\title{Résumé du cours de Methodologie}
\author{Salomon Heya}
\date{March 2026}

\begin{document}

\maketitle

\section{Objectifs principaux et Adaptabilité}

Le PIU (Processus Investigatif Universel) comprend plusieurs objectifs fondamentaux à savoir: 
\begin{itemize}
    \item L'Etablissement factuel objectif qui consiste à la documentation des faits de manière reproductible et vérifiable
    \item L'Intégrité probatoire qui consiste à garantir que les éléments collectés conservent leur valeur probante
    \item La Traçabilité complète qui vise à permettre l'audit de chaque étape du processus investigatif
    \item La Défendabilité, il doit produire des conclusions résistant à la contestation contradictoire
\end{itemize}

Le PIU possède un principe d'adaptabilité qui permet son application dans 3 contextes :
\begin{itemize}
    \item Contexte Particulier
    \item Contexte Organisationnel
    \item Contexte Judiciaire
\end{itemize}

\section{Les Quality Gates}
Les Quality Gates constituent des points de contrôle obligatoires validant la conformité avant proression vers la phase suivante, leur structure s'articule autour de critères graduels (Obligatoires, standards, avancés)
Les différentes phases des Quality Gates sont:
\begin{itemize}
    \item Phase 1 : Initialisation et Cadrage qui determine la récevabilité de la demande investigative et établit les conditions opérationnelles de sa conduite. Elle comprend plusieurs étapes;
    \begin{itemize}
        \item Etape 1 : Réception et Evaluation primaire
        \item Etape 2 : Décision d'engagement
        \item Etape 3 : Planification Opérationnelle  
    \end{itemize} 
    \item Phase 2 :Acquisition et Préservation qui vise à constituer la base probatoire de l'investigation en garantissant l'intégrité et la recevabilité juridique des éléments collectés
    \begin{itemize}
        \item Etape 4 : Collecte Systématique des Elements probatoires
        \item Etape 5 : Sécurisation et Préservation
        \item Etape 6 : Identification et Cartographie des Parties Prenantes
    \end{itemize}
    \item Phase 3 : Invstigation Analytique qui constitue le coeur du processus; elle comporte une boucle itérative formalisée
\begin{itemize}
    \item Etape 7 : Interrogatoire et Collecte testimoniale
    \item Etape 8 : Génération d'Hypothèses Multiples
    \item Etape 9 : Confrontation Hypothèses- Preuves et la Falsification Active
    \item Etape 10 :Validation et Vérification
\end{itemize}

\item Phase 4 : Formalisation et Transmission qui consiste à transformer le travail investigatif en livrables structurés, adaptés au destinataire et juridiquement recevables
\begin{itemize}
    \item Etape 11 : Documentation Exhaustive et Consolidation
    \item Etape 12 : Rédaction du Livrable Final 
    \item Etape 13 : Transmission sécurisée du livrable
\end{itemize}
    \item Phase 5 : Suivi Post-Transmission : Une phase cruciale dans le cas des investigations judiciaires ou de cintestations anticipées
    \begin{itemize}
        \item Etape 14 : Préparation à la Défense Contradictoire
        \item Etape 15 : Maintenance Probatoire Continue
        \item Etape 16 : Analyse des Contre-Expertises
        \item Etape 17 :Clôture et retour d'experience
    \end{itemize}
\end{itemize}
Les différents Qualities Gates sont : 
\begin{itemize}
    \item Quality Gate 1 : Validation Pré-Opérationnelle qui vise à garantir que l'ensemble des conditions préalables au démarrage effectif de l'investigation sont satisfaites
    \item Quality Gate 2 : Intégrité Probatoire, il vise à vérifier que la base probatoire constituée est complète, intègre, traçable et juridiquement exploitable avant engagement de la phase analytique
    \item Quality Gate 3 : Robustesse Analytique : on cherhce ici à s'assurer que les conclusions investigatives sont suffisamment robustes, documentées et defendables pour justifier leur formalisation et transmission au commanditaire
\end{itemize}


\end{document}
